\documentclass[10pt,a4paper]{amsart}
\usepackage[margin=19mm]{geometry}
\usepackage{amsmath,amssymb,mathtools}
\usepackage[scr=rsfs,cal=euler]{mathalfa}
\usepackage{xfrac}
\usepackage[unicode,hidelinks,bookmarks=false]{hyperref}
\newcommand{\ie}{i.e.\ }
\newcommand{\cf}{cf.\ }
\newcommand{\resp}{respectively\ }
\newcommand{\Subsection}{\subsection}
\providecommand{\nobreakdash}{\nobreak\hbox{-}}
\setlength{\emergencystretch}{2em}
\multlinegap0pt
\usepackage{multirow}
\usepackage{amssymb}
\usepackage{mathtools}
\usepackage{xcolor}
\usepackage{bbm}
\newtheorem{lem}{Lemma}
\usepackage{cleveref}
\crefname{lem}{Lemma}{Lemmas}
\newcommand{\GL}{\rm{GL}}
\newcommand{\R}{\bb R}
\newcommand{\Z}{\bb Z}
\newcommand{\bb}[1]{\mathbb{#1}}
\renewcommand{\rm}[1]{\mathrm{#1}}
\title{Density of shapes of periodic tori in the cubic case}
\author{Nguyen-Thi Dang, Nihar Gargava, Jialun Li}
\date{Source: arXiv:2502.12754v2 (CC BY 4.0)}
\begin{document}
\maketitle

\noindent\emph{Source and licence.} Density of shapes of periodic tori in the cubic case. Version arXiv:2502.12754v2 (CC BY 4.0), distributed by arXiv under the Creative Commons Attribution 4.0 International licence, http://creativecommons.org/licenses/by/4.0/, as recorded in the arXiv licence metadata for this exact version and shown on the abstract page https://arxiv.org/abs/2502.12754v2. Full licence notice: \texttt{licenses/arxiv-2502.12754-CC-BY-4.0.txt}.\par\smallskip

\noindent\textit{Editorial scope.} Counted unit: the article's lem defin-Lcde (source0.tex lines 1171--1198). The article's complete original proof of it is delivered with it (source0.tex lines 1199--1318, one \texttt{proof} environment of 120 lines). No mathematics is written, repaired or re-derived: the statement and the proof are the article's own lines, in the article's own notation, and no retained line is substituted or re-flowed.  Retained uncounted as the source-local context the proof needs: lines 1155--1166.  Nothing was omitted from any retained span, so the package carries no omission record.  No retained line contains a citation, so the wrapper carries no bibliography and no bibliography entry.  No retained line contains a figure environment, an image-inclusion command or drawing code, so no image is delivered and the package declares no asset dependency.  Editorial formatting only, in addition: an AMS \texttt{amsart} preamble that restates the article's own declarations and macros used by the retained lines, namely the environments \texttt{lem} on separate counters, together with the standard packages those lines need.  The retained environments print in this document as Lemma 1 in order of appearance, whereas the article prints the same items with the numbering of its own full document.  The complete native source of the article as distributed in the arXiv e-print for this exact version is bundled unchanged as \texttt{sources/173469/source0.tex}.
	\begin{equation}\label{eq:defi_of_L}
		L_{c,d,r}=
		\left\lbrace\begin{tabular}{c|ccl}
\multirow{3}*{$\begin{pmatrix}
    m \\
    n
\end{pmatrix} \in \Z^2$} 
& $m+n$  &$\equiv\,0$ &  $\pmod {7\cdot 2^{c-1}}$ \\ 
& $m-2n$ &$\equiv\,0$ &  $\pmod {8\cdot 3^{d-1}}$ \\ 
& $2m-n$ &$\equiv\,0$ &  $\pmod {24\cdot 5^{r-1}}$ 
		\end{tabular}\right\rbrace
	\end{equation}

	\begin{lem}\label{defin-Lcde}
		Denote by 
		$$g = 8 
		\left( 
		\begin{array}{rr}
			1 & 1\\
			0 & -1  
		\end{array}\right)
		\begin{pmatrix}
			3 & 0 \\
			0 & 1
		\end{pmatrix}
		\left( \begin{array}{rr}
			-1 & 1 \\
			2 & -1
		\end{array}\right).$$
		Then for every $c\geq 4$ and $d\geq 2$ and $r\geq 1$, the sub-lattice $L_{c,d,r}$ as defined in \cref{eq:defi_of_L} satisfies 
		$$L_{c,d,r} = g
		\begin{pmatrix}
			3^{d-2} & 0 \\
			0 & 5^{r-1}
		\end{pmatrix} 
		\bigg\lbrace 
		\binom{x}{y} \in \Z^2
		\; \bigg\vert \;
		3^{d-2} x - 5^{r-1} y \equiv 0 \pmod{ 7 \cdot 2^{c-4}}
		\bigg\rbrace.$$
	\end{lem}

	\begin{proof}
		By \cref{eq:defi_of_L}, $L_{c,d,r}$ is the set of integer vectors $(m,n) \in \Z^2$ that solve a rank-two equation in $\Z^3$, i.e., such that 
		$$ 
		\left(\begin{array}{rr}
			1 & 1 \\
			1 & -2 \\
			2 & -1
		\end{array}\right)
		\binom{m}{n} \in 
		\left(\begin{array}{ccc}
			7 \cdot 2^{c-1} & & (0) \\
			& 8 \cdot 3^{d-1} & \\
			(0) & & 24 \cdot 5^{r-1}
		\end{array}\right) \Z^3 .$$
		We compute the Smith normal form on the $3 \times 2$ matrix on the left and denote by 
		$$ 
		U = \left( 
		\begin{array}{rrr}
			0 &  -1 & 1\\
			0 & 2 &  -1\\
			1 & 1 &  -1
		\end{array}\right) , \quad
		M = 
		\left(
		\begin{array}{rr}
			1 & 1 \\
			1 & -2 \\
			2 & -1 
		\end{array}\right), \quad
		V = \left( 
		\begin{array}{rr}
			1 & 1\\
			0 & -1  
		\end{array}\right).
		$$ 
		Then $U \in \GL(3,\Z)$ and $V \in \GL(2,\Z)$ and
		$     M = U^{-1} \left(
		\begin{array}{rr}
			1 & 0 \\
			0 & 3 \\
			0 & 0
		\end{array} \right) V^{-1}.
		$
		Now, $V^{-1} L_{c,d,r}$ is the set of $(m,n)\in \Z^2$ such that there exists $(x,y,z) \in \Z^3$ such that
		\begin{equation}\label{eq:nicer-Lcde}
			\left(
			\begin{array}{rr}
				1 & 0 \\
				0 & 3 \\
				0 & 0
			\end{array} \right) 
			\binom{m}{n} =
			U  \left(\begin{array}{r}
				7 \cdot 2^{c-1} \; x  \\
				8 \cdot 3^{d-1} \; y \\
				24 \cdot 5^{r-1}\; z
			\end{array}\right)  . 
		\end{equation}
		The first two lines of the matrix equation \cref{eq:nicer-Lcde} yield the following rank-$2$ linear equation
		\begin{equation}\label{eq:Lcde-2square}
			\left( 
			\begin{array}{rr}
				1 & 0\\
				0 & 3
			\end{array}\right)
			\binom{m}{n} =
			\left( 
			\begin{array}{rr}
				-1 & 1\\
				2 &  -1
			\end{array}\right)
			\binom{8 \cdot 3^{d-1} \; y}{24 \cdot 5^{r-1}\; z}. 
		\end{equation}
		Hence using that $d \geq 2$, we deduce that 
		$V^{-1} L_{c,d,r} \subset 
		8 \begin{pmatrix}
			3 & 0 \\
			0 & 1
		\end{pmatrix} 
		\left( 
		\begin{array}{rr}
			-1 & 1\\
			2 &  -1
		\end{array}\right)
		\begin{pmatrix}
			3^{d-2} & 0\\
			0 &  5^{r-1}
		\end{pmatrix} \Z^2.
		$ 
		
		The last line of the matrix equation \cref{eq:nicer-Lcde} is a hyperplane equation in $\Z^3$. 
		It will induce a sub-lattice in $\Z^2$ which will be mapped to $L_{c,d,r}$ by an element in $\GL(2,\R)$.
		By the matrix equation, $(x,y,z)\in \Z^3$ satisfy
		\begin{equation}\label{eq:Lcde-hyperplane}
			7 \cdot 2^{c-1} \; x + 8 \cdot 3^{d-1} \; y - 24 \cdot 5^{r-1}\; z = 0.
		\end{equation}
		Equivalently,
		$$ -7 \cdot 2^{c-4} \; x = 3^{d-1} \; y - 3 \cdot 5^{r-1}\; z. $$
		Since $c \geq 4$ and using that $3$ and $7 \cdot 2^{c-4}$ are coprime and $d \geq 2$, the integers $(y,z)\in \Z^2$ solve the congruence equation
		$ 3^{d-2} \; y -  5^{r-1}\; z \equiv 0 \; \pmod{ 7 \cdot 2^{c-4}}.$
		
		Finally, combining the hyperplane relation with the rank $2$ part, we deduce that
		$$ L_{c,d,r} = 8 V \begin{pmatrix}
			3 & 0 \\
			0 & 1
		\end{pmatrix} 
		\left( 
		\begin{array}{rr}
			-1 & 1\\
			2 &  -1
		\end{array}\right)
		\begin{pmatrix}
			3^{d-2} & 0\\
			0 &  5^{r-1}
		\end{pmatrix} \bigg\lbrace 
		\binom{y}{z}
		\; \bigg\vert \;
		3^{d-2} y - 5^{r-1} z \equiv 0 \pmod{ 7 \cdot 2^{c-4}}
		\bigg\rbrace. $$
	\end{proof}

\end{document}
